\subsection{Metadata Schema and Field Descriptions}
\label{app:prompts_schema}

Each prompt in the bank conforms to a standardized JSON schema validated at harness startup (\cref{subsubsec:config_validation}). The schema enforces type safety, required fields, and enumerated category values to prevent malformed inputs from biasing results. Field definitions are as follows:

\begin{itemize}[leftmargin=*]
    \item \texttt{id} (string, required): A unique identifier in the format \texttt{<category\_prefix>\_<number>} (e.g., \texttt{soc\_001}, \texttt{ctx\_012}). IDs enable unambiguous cross-referencing between raw model outputs (\cref{subsubsec:data_persistence}), human rubric scores (\cref{app:socratic_quality_rubric,app:context_fidelity_rubric}), and statistical analyses (\cref{sec:statistical_analysis}). The prefix encodes category membership for stratified sampling and category-specific subgroup analyses (\cref{subsec:comparative_analysis}).

    \item \texttt{category} (string, required): One of \texttt{["Socratic Probing", "Adaptive Contextualization", "Feedback \& Improvement", "Edge Cases \& Adversarial", "Safety Challenges", "Physical Estimations", "Order-of-Magnitude Checks", "Constraints \& Trade-offs", "Error Detection \& Correction", "Multi-turn Dialogue", "Metacognitive Prompting", "Domain Transfer"]}. This field corresponds to the Intent Category taxonomy defined in \cref{app:prompts_taxonomy} and drives the harness's category-based execution modes (\cref{subsubsec:execution_modes}) and enables filtering for targeted experiments (e.g., context-only runs; \cref{subsec:contextualization_test}).

    \item \texttt{prompt\_text} (string, required): The exact student utterance presented to the model, including any simulated errors, questions, or metacognitive statements. Text is prefixed with \texttt{"Student: "} to clearly demarcate role boundaries and reinforce the Socratic tutoring context established in the system prompt (\cref{subsec:datasets_prompt_bank}). All prompts are self-contained within a single turn; multi-turn exchanges were simulated by sequential prompt issuance in contextualization experiments (\cref{subsec:contextualization_test}).

    \item \texttt{teacher\_context} (string or null, required): Instructor-provided, problem-specific information dynamically supplied at runtime to simulate Mettle's adaptive contextualization feature. When non-null, this field contains numerical constraints, physical parameters, system specifications, or domain-specific hints that the model should incorporate into its response. Context strings are injected into the user message (ephemeral strategy; \cref{subsec:contextualization_test}) or system message (persistent strategy) depending on experimental configuration. Null values indicate prompts designed to test general reasoning without additional context (typical for Socratic Probing and Safety Challenges categories).

    \item \texttt{expected\_keywords} (array of strings, required): A curated list of terms, phrases, or concepts that a high-quality response should address. Keywords serve dual purposes: (a) as ground truth for automated context fidelity checks (\cref{subsec:automated_metrics}), where keyword inclusion rates are computed via case-insensitive substring matching, and (b) as qualitative guidance for human raters when applying the Socratic and Context Fidelity rubrics (\cref{sec:scoring_rubrics}). Keywords were populated through expert review and refined during pilot rating sessions to ensure inter-rater reliability (IRR $\kappa > 0.4$; \cref{subsec:sampling_blinding}).
\end{itemize}

Additional metadata fields not shown in the JSON but logged during execution include: \texttt{system\_prompt\_template} (identifier for the Socratic/neutral/hybrid variant used; \cref{subsec:few_shot_prompts}), \texttt{temperature} (sampling parameter; \cref{subsec:temperature_sweep}), \texttt{timestamp\_utc} (execution time for temporal analysis), \texttt{model\_name} and \texttt{model\_version} (for version tracking; \cref{subsubsec:adapter_interface}), and \texttt{run\_id} (UUID for multi-run replication; \cref{subsec:determinism}).
